\section{Abstract Framework}
\label{section:44-GAI-Framework}

The \gai benchmark framework is presented below and summarized in figure \osullikomment{TBC}. 
The framework is the overarching term given to the collection of \gai problems that comprise the benchmark. 
Each problem is separately designed and implemented in the benchmark framework by the methodology described in section~\ref{section:44-GAI-Design}.
This section will present each step again, highlighting the \textit{general} design choices at this framework level. 
Chapter~\ref{chapter:50-Community-Detection-Case-Study} demonstrates a worked example of how a benchmark problem is implemented. 


\osullikomment{Insert updated version of variance table here with the constraints added as the figure here.}

\textbf{Identifing and Characterizing a Relevant Problem.}
The ARLIS domain study~\cite{Regli2022a} and the extended problem survey presented in table~\ref{table:problem-survey}, identifies $11$ \gai problems relevant to the domain of interest. 
If a problem not from the list in table~\ref{table:problem-survey} is selected, evidence of its importance must also be documented. 
Further, the problems should be defined as \textit{problems}, not specific \textit{kernels} to execute. 
For example, a problem definition might be \textit{Given the input X, identify who the immediate associates of a person with name Y, provide the output as a list of vertex IDs}. 
Rather than saying \textit{Run BFS from node Y} the problem is reflective of the actual domain, where an analyst knows the name (a vertex attribute) of a target, but not their vertex ID and there is no specification about how to solve it.
One implementation might load the entire graph, BFS for the target node and then BFS for their neighbors. 
Another might be to load the entire graph, DFS for the target node and then BFS for the neighbors. 
Another might load the entire graph into a GDB and use an index lookup and yet another might examine each vertex as it is created and take a chance on early stopping by pausing graph creation $1000$ vertices after the target is identified, hoping that the complete neighbor list is created by then. 
The problem-based approach still allows for the comparison of like-approaches, but also offers the flexibility of solving problems in other ways, allowing someone evaluating \gai systems using the benchmark to test if using a particular approach on one problem that may be slower may be offset by an improvement to another system. 

\textbf{Identifing and Characterizing the Phenomena of Interest.}
The phenomena of interest is the trade-off between the efficiency, effectiveness and cost of \gai systems.
Thus, the metrics we expect to be relevant to all \gai tasks are \textit{ETL Time, Execution Time, Batch Latency\footnote{\textnormal{\scriptsize The time between when a vertex or edge is added to a graph and when it appears in the result of a workload~\cite{Basak2020}}}} and \textit{Cost}. 
For each specific problem, the phenomena must be extended to articulate the specific effectiveness metrics, so that the assessed `goodness' of an output is meaningful.

\textbf{Conduct Exploratory Experimentation.}
For the \gai Benchmark, we review the results of workload analyses to identify the bottlenecks in systems and their causes.
We identify that for graph processing, the input data size, shape and velocity impact the efficiency and cost of benchmarks, and that the flow-on effects of memory, computational and communication bottlenecks are dependent on the input data. 
All \gai benchmarking problems must therefore have at least one input graph that is in the scale-free, uniform-random and mesh categories. 
Preferably, these will be available at tiny (processing is possible entirely in the cache of a commercial system, $<60$Mb\footnote{\scriptsize{Based on Intel Xeon Gold \url{https://www.intel.com/content/www/us/en/products/details/processors/xeon/scalable/gold/products.html}}}), small (processing is possible in memory on a personal computer, no larger than $\sim16$Gb), medium (processing is possible on a commercial system, no larger than $\sim1$Tb) and large (processing will require multiple commercial systems (at least $\sim5$Tb)).

\textbf{Deveop Hypotheses.}
For the \gai benchmark, the hypotheses should use the three constraint levers of time (efficiency), resources (cost) and accuracy (effectiveness) to design experiments. 
These levers provide a mechanism to identify the trade-offs in \gai between efficiency, effectiveness and cost by constraining one (or more) and observing changes in the phenomena under observation.
For fairness, each \gai problem should allow an unconstrained run before imposing constraints in subsequent iterations.
Table~\ref{table:gai-variance} summarizes the hypotheses expected to hold across all problems in the \gai benchmark.

% \begin{landscape}
    
    \begin{table*}[t]
        \begin{center}
        \small
            \begin{tabular}{ll|cccl} 
                Dimension                   &   Components              & Effectiveness       & Efficiency      &  Cost         &   Notes         \\
                \hline
                \textbf{Specification}      &   Datasets                & $\checkmark$        & $\checkmark$    & $\checkmark$  & Size and Shape            \\
                                            &   Tasks                   & $\checkmark$        & $\checkmark$    & $\checkmark$  & In Aggregate              \\
                                            &   Metrics                 &                     &                 &               & Human Perception          \\
                            ~               &           ~               &         ~           &         ~       &       ~       &                           \\
                \textbf{System under Test}  &   Hardware                &                     & $\checkmark$    & $\checkmark$  & Capital \& Ops cost \\
                                            &   Software                &                     & $\checkmark$    & $\checkmark$  & Capital \& Ops cost \\
                                            &   Implementation          & $\checkmark$        & $\checkmark$    & $\checkmark$  &                           \\
                            ~               &           ~               &         ~           &         ~       &       ~       &                           \\
                \textbf{Constraints}        &   None                    &                     &                 &               & Baseline                  \\
                                            &   Time                    & $\downarrow$        &                 & $\downarrow$  &                           \\
                                            &   Accuracy                &                     & $\uparrow$      & $\downarrow$  &                           \\
                                            &   Cost                    & $\downarrow$        & $\downarrow$    &               &                           \\

            \end{tabular}
        \normalsize
        \caption{Summary of general hypotheses for the \gai benchmark} 
        \label{table:gai-variance}
        \end{center}
    \end{table*}
% \end{landscape}

\noindent Hypotheses should be formally specified in the form:

{\centering \textit{When [lever] is constrained for [task] on [dataset] we should observe a [description of expected change] in [phenmonena]}.}

\textbf{Construct the Investigation Apparatus.}
For the \gai benchmark, we identify the need to track \textit{time} as a proxy for efficiency, \textit{accuracy} (or equivalent) as a proxy for effectiveness and \textit{power} as a proxy for cost.
In a purely Input/Output paradigm with access only to the external view of the system under test, time is measured from when the input data is sent until the output is received. 
More detailed timing data generally requires feedback from a system. 
It can either time itself and report timings or needs to send proof of progress as execution proceeds.
To promote the exploration of trade-offs between efficiency and effectiveness, we propose a 'polling' system that will periodically ask the system under test for a result, and continue to do so until the system has `completed' its workload.
The results of the `polling' will reveal how accuracy covaries with time. 

The accuracy is calculated by examining the provided answer against a reference answer.
Instinctively, we see exact solutions as required for many \gai problems, particularly those in the SPS problem space, but if the problem is framed correctly, they can accept approximate solutions. 
For example, \textit{Given the input X, identify who the immediate associates of a person with name Y, provide the output as a list of vertex IDs} will have an exact answer, $N$, but there is still utility in the real world for an answer that is $N\pm1$ or in a scale-free graph maybe even $N\pm100$.
Accepting a distance from the exact solution here is now an interesting metric that allows us to perhaps trade off the accuracy of a solution for a faster solution, that is much quicker. 

In the Input/Output paradigm, there is no easy way to measure power (cost).
One method is to estimate it using some kind of cost model that would require the SUT owner to enter information about the system that can project power usage over time. 
Another is to provide a hosted benchmark and use that to establish a baseline of power consumption. 
A third is to provide a microbenchmark and instrumentation instructions so that the SUT user can generate power consumption statistics separately and report them. 
None of these approaches are optimal, so for the \gai benchmark, we explore using a hosted instance of the benchmark to measure power consumption but focus efforts on the tradeoff between efficiency and effectiveness in the first instance. 

\osullikomment{Fill in QUICC system details here.}

\textbf{Analyze the Results, Developing a Theoretic Explanation and Iterating.}
The \gai benchmark should report the \textit{efficiency, effectiveness} and \textit{cost} of each workload on each system under test configuration. 
The results should be grouped in one dimension by the system under test, showing its relative performance across all the \gai benchmark problems. 
They should also group by problem so that the performance of different systems under test on the same problems can be compared. 

\begin{figure}[h]
    \centering
    \newcommand{\radarplot}[6]{ % Define a command for radar plot
        \begin{tikzpicture}
            \begin{polaraxis}[
                ymax=1.1, ymin=0, 
                xtick=\empty, ytick=\empty,
                width=0.3\linewidth,
                title={#1},
                ylabel={#2},
                xticklabels={E,C,A},
            ]
                \addplot+[mark=none,smooth cycle,fill=#6] coordinates {
                    (0,0) (90,#3) (210,#4) (330,#5)
                };
            \end{polaraxis}
        \end{tikzpicture}
    }

    \begin{tabular}{cccc}
             & Problem 1                                & Problem2                              & Problem3                              \\
        SUT1 & \radarplot{}{}{1}{1}{1}{blue!30} & \radarplot{}{}{0.3}{0.7}{0.9}{red!30} & \radarplot{}{}{0.8}{0.5}{0.6}{green!30} \\
        SUT2 & \radarplot{}{}{0.7}{0.9}{0.4}{orange!30} & \radarplot{}{}{0.5}{0.5}{0.5}{purple!30} & \radarplot{}{}{0.4}{0.9}{0.7}{yellow!30} \\
        SUT3 & \radarplot{}{}{0.8}{0.1}{0.9}{cyan!30} & \radarplot{}{}{0.6}{0.5}{0.7}{magenta!30} & \radarplot{}{}{0.9}{0.7}{1.0}{teal!30} \\
    \end{tabular}

    \caption{Mock-Up of Radar Plots to show what the output of a \gai benchmark run looks like. Radar plots are along the axes of effectiveness (north) efficiency (southwest) and cost (southeast)}
    \label{fig:radar_plots}
\end{figure}



% \begin{landscape}
    
    \begin{figure}[h!]
        \begin{center}
            \includesvg[width=\columnwidth]{Figures/SVGs/benchmark_design_process}
            \caption{\gai Benchmark Design Framework}
            \label{figure:framework}
        \end{center}
    \end{figure}

% \end{landscape}